\chapter{Fractional Gaussian Fields}\label{chap:gaussianFields}
In this chapter we will be constructing the fractional Gaussian fields, that we will later see converge in distribution. The construction of the fractional Gaussian fields very much depends on the Riesz kernel and the Dirichlet Laplacian.

We will also give an example of an discrete fractional Gaussian field, which we will later see in some simulations.

\section{Discrete Fractional Gaussian Fields}

\begin{defi}\label{fractionalGaussianFieldDiscrete}
	Let $s \geq 0$, then a \emph{discrete fractional Gaussian field}$ X_s^{m}$ with parameter $s$ on $V_m$ is a Gaussian vector indexed by $V_m$ with mean zero and covariance matrix $\mathbf{\Sigma} = \left(G_{2s}^{m}(x,y)\right)_{x,y \in V_m}$
\end{defi}

For any $f \in \ell(V_m)$ we will use the notation
\begin{equation}
	X_s^{m}(f) = \frac{1}{a_m}\sum_{p \in V_m}f(p)X^{m}_{s}(p)
\end{equation}


We then note that the discrete fractional Gaussian field has the following property
\begin{prop}\label{theOneWeNeedInTheEndSomehow}
	For $f,g \in \ell(V_m)$ we have that
	\begin{equation}
		\e \left[X_s^{m}(f)X_s^{m}(g)\right] = \frac{1}{a_m} \sum_{p \in V_m} (-\Delta_{m})^{-s}f(p)(-\Delta_{m})^{-s}g(p)
	\end{equation}
\end{prop}
\begin{proof}
	Using \cref{discreteFractionalLaplacianProperties} we find
	\begin{align}
		\e \left[X_s^{m}(f)X_s^{m}(g)\right] &= \frac{1}{2a_m} \sum_{p,q \in V_m}f(p)g(q)\e\left[X_s^{m}(p)X_s^{m}(q)\right] \\
		&=  \frac{1}{2a_m} \sum_{p,q \in V_m}f(p)g(q)G_{2s}^{m}(p,q) \\
		&= \frac{1}{a_m} \sum_{p \in V_m}f(p)\left(\frac{1}{a_m} \sum_{q \in V_m}g(q)G_{2s}^{m}(p,q)\right) \\
		& = \frac{1}{a_m} \sum_{p \in V_m}f(p)(-\Delta_{m})^{-2s}g(p) \\
		& = \frac{1}{a_m} \sum_{p \in V_m} (-\Delta_{m})^{-s}f(p)(-\Delta_{m})^{-s}g(p)
	\end{align}
\end{proof}

\begin{example}\label{OneToSimulate}
	Here we will give an example of a discrete fractional Gaussian field. Let $(W_i)_{1\leq i \leq N_m}$ be a sequence of i.i.d Gaussian random variables with mean 0 and variance 1, i.e. white noise. Then define
	\begin{equation}
		X_{s}^{m}(x) = \sum_{i=1}^{N_m}\left(\lambda_i^{m}\right)^{-s}\phi^{m}_{i}(x)W_i, \quad x \in V_m.
	\end{equation}
	Initially, we check the mean, so
	\begin{equation}
		\e \left[X_s^{m}\right] = \begin{pmatrix}
			\e\left[X_s^{m}(x_1)\right] \\
			\vdots \\
			\e \left[X_s^{m}(x_{N_m})\right]
		\end{pmatrix} = \begin{pmatrix}
		0 \\
		\vdots \\
		0
		\end{pmatrix}
	\end{equation}
	since
	\begin{equation}
		\e \left[X_s^{m}(x_i)\right] = \e \left[\sum_{i=1}^{N_m}\left(\lambda_i^{m}\right)^{-s}\phi^{m}_{i}(x)W_i\right] = \sum_{i = 1}^{N_m} \left(\lambda_i^{m}\right)^{-s}\phi^{m}_{i}(x) \e \left[W_i\right] = 0.
	\end{equation}
	Moving on to the covariance matrix, we have that
	\begin{multline}
		\mathbf{\Sigma} = \e \left( \left(X_s^{m} - \e \left[X_s^{m}\right] \right)\left(\left(X_s^{m}\right)^{T} - \e \left[X_s^{m}\right] \right)\right) = \\
		\e \left[X_s^{m} \cdot \left(X_s^{m}\right)^{T}\right] = \e \left[\begin{pmatrix}
			X_s^{m}(x_1)X_s^{m}(x_1) & \cdots  & X_s^{m}(x_1)X_s^{m}(x_{N_m}) \\
			\vdots & \ddots  & \vdots \\
			X_s^{m}(x_{N_m})X_s^{m}(x_1)  & \cdots & X_s^{m}(x_{N_m})X_s^{m}(x_{N_m}) 
		\end{pmatrix}\right].
	\end{multline}
	Looking at a single element $\mathbf{\Sigma}_{p,q}$ we find 
	\begin{align}
		\e \left[X_{s}^{m}(x_p)X_s^{m}(x_q)\right] &= \e \left[\left(\sum_{i=1}^{N_m}\left(\lambda_i^{m}\right)^{-s}\phi^{m}_{i}(x_p)W_i\right) \left(\sum_{j=1}^{N_m}\left(\lambda_i^{m}\right)^{-s}\phi^{m}_{i}(x_q)W_i\right)\right]\\ 
		&= \e \left[\sum_{i=1}^{N_m}\left(\lambda_i^{m}\right)^{-2s}\phi^{m}_{i}(x_p)\phi^{m}_{i}(x_q)W_i^2\right] +\e \left[\sum_{i,j=1, i \neq j }^{N_m}\left(\lambda_i^{m}\lambda_j^{m}\right)^{-s}\phi^{m}_{i}(x_p) \phi^{m}_{j}(x_q)W_iW_j\right] \\
		&= \sum_{i=1}^{N_m}\left(\lambda_i^{m}\right)^{-2s}\phi^{m}_{i}(x_p)\phi^{m}_{i}(x_q)\e \left[W_i^2\right] +\sum_{i,j=1, i \neq j }^{N_m}\left(\lambda_i^{m}\lambda_j^{m}\right)^{-s}\phi^{m}_{i}(x_p) \phi^{m}_{j}(x_q)\e \left[W_i\right] \e \left[W_j\right] \\
		&= \sum_{i=1}^{N_m}\left(\lambda_i^{m}\right)^{-2s}\phi^{m}_{i}(x_p)\phi^{m}_{i}(x_q) \\
		&= G_{2s}^{m}(x_p,x_q)
	\end{align}
	using that $W_i$ are independent. Thus $X_s^{m}$ is a discrete fractional Gaussian field.
\end{example}

\section{Fractional Gaussian Fields}

\begin{defi}\label{fracGaussianField}
	Let $s \geq 0$. A \emph{fractional Gaussian field} $X_s$ with parameter $s$ on $K$ is a centered Gaussian field $\cbr{X_s(f): f \in S(K)}$ such that for any $f,g \in S(K)$
	\begin{equation}
		\e \left[X_s(f) X_s(g)\right] = \int_{K} (-\Delta)^{-s}f (-\Delta)^{-s}g \dd \mu
	\end{equation}
\end{defi}

We now move on to proving that the fractional Gaussian fields exist. First we will state the Bochner-Minlos theorem for nuclear spaces, since it is needed for the proof of existence. It will be given without proof but one can be found in \cite{minlos}.

\begin{theorem}[Bochner-Minlos]\label{BochnerMinlos}
	Assume $\mathcal{X}$ is a nuclear space. Any continuos positive-definite linear functional $\Lambda$ on $\mathcal{X}$ satisfying $\Lambda(0) = 1$ is the Fourier transform of a countably additive positive normalized measure on $\mathcal{X}'$
\end{theorem}

We will first state a lemma, to prove that a functional, that we will need, is positive definite.
\begin{lemma}\label{LambdaPosDefinite}
	The functional $\Phi: \mathcal{X} \to \ree$ defined as  $\Phi(v) = \exp(-\frac{1}{2} \inn{v}{v})$ is positive definite.
\end{lemma}
\begin{proof}
	Let $v_1, \ldots v_n$ be elements of $\mathcal{X}$ and choose an orthonormal basis $e_1, \ldots e_m$ of $\text{span}\cbr{v_1, \ldots, v_n}$. Further let $Z = (Z_1, \ldots, Z_m)$ be a vector of independent standard normal random variables, and note that for all $u \in \ree^{m}$ we have that
	\begin{equation}
		\Phi \left(\sum_{j = 1}^{m} u_j e_j\right) = \exp \left(-\frac{1}{2} \sum_{i = 1}^{m} u_i^2\right) = \e \left[e^{ \img u \cd Z}\right].
	\end{equation}
	This leads us to for $z \in \cee^{n}$ 
	\begin{equation}
		\sum_{j, k = 1}^{n}z_j \overline{z_k} \Phi(v_j - v_k) = \sum_{j,  k = 1} z_j \overline{z_k} \e \left[e^{\img (v_j -v_k) \cd Z}\right] = \e \left[\abs{\sum_{j = 1} z_j e^{i v_j \cdot Z}}^2\right] \geq 0.
	\end{equation}
\end{proof}

\begin{theorem}
	Let $s \geq 0$. There exists a centered Gaussian distribution $X_s$ on $S'(K)$ such that for $f,g \in S(K)$
	\begin{equation}
		\e[X_s(f)X_s(g)] = \int_{K} (-\Delta)^{-s}f (-\Delta)^{-s}g \dd \mu
	\end{equation}
\end{theorem}

\begin{proof}
	The space $S(K)$ is a nuclear space, why it is enough to use \cref{BochnerMinlos} (Bochner-Minlos) on the functional
	\begin{equation}
		\Lambda : f \to \exp\left(-\frac{1}{2} \norm{f}^2_{k}\right).
	\end{equation}
	Since if \cref{BochnerMinlos} hold, then it gives us an unique probability measure $\nu$ such
	\begin{equation}
		\exp \left(-\frac{1}{2} \norm{f}^{2}_{k}\right) = \int_{S'(K)} e^{i \inn{x}{f}} \dd \nu (x).
	\end{equation}
	Notice that the right side is a characteristic function for some random variable. If we choose $f= t\cd g$ for some $g \in S(K)$ and $t \in \ree$ (note that $t \cd g \in S(K)$) then the left side of the equation becomes the characteristic function for the normal distribution with mean zero and variance $\norm{f}_{s}^{2}$. Thus we have the $X_s$ we wanted.
	
	
	Now we prove that $\Lambda$ satisfies the requirements of \cref{BochnerMinlos}. Obviously $\Lambda(0) = 1$, and from \cref{LambdaPosDefinite} $\Lambda$ is positive definite. Lastly we need to prove it is continuous at 0. 
	
	First we note that since $(-\Delta)^{-s}$ is a bounded operator we have that there exists some constant $C> 0$ so $\norm{(-\Delta)^{-s}f }_{L^{2}(K, \mu)} \leq C \norm{f}_{L^{2}(K, \mu)}$. Now let $\epsilon > 0$ be given, choose $$\delta = \begin{cases}
		\left(\frac{-2\cd \ln(1-\epsilon)}{C}\right)^{\frac{1}{2}}, & \text{when } \epsilon \in (0,1) \\
		1,& \text{when } \epsilon \geq 1.
	\end{cases}$$
	First if $\epsilon \geq 1$ notice that 
	\begin{equation}
		\abs{\exp\left(-\frac{1}{2} \norm{f}_{s}^{2}\right) - \exp\left(-\frac{1}{2} \norm{0}_{s}^{2}\right)} \leq 1 \leq \epsilon, \quad \forall f \in S(K)
	\end{equation} 
	in which case we can choose $\delta$ to be any real number larger than $0$. Moving on to $\epsilon \in (0,1)$ we get
	\begin{align}
		\abs{\exp\left(-\frac{1}{2} \norm{f}_{s}^{2}\right) - \exp\left(-\frac{1}{2} \norm{0}_{s}^{2}\right)} &\leq \abs{\exp\left(-\frac{1}{2} C\norm{f}_{L^2(K,\mu)}\right) - 1} \\
		&= \abs{\exp\left(-\frac{1}{2} C \int_{K}\abs{f}^2 \dd \mu\right) - 1} \\
		&\leq \abs{\exp\left(-\frac{1}{2} C \int_{K} \delta^2 \dd \mu\right) - 1} \\
		&= \abs{\exp\left(-\frac{1}{2} C \int_{K} \frac{-2\cd \ln(1-\epsilon)}{ C} \dd \mu\right) - 1}  \\
		&= \abs{\exp\left(\ln(1-\epsilon)\right) -1}\\
		& = \epsilon
	\end{align}
	using that $\mu$ is a probability measure in the third to last equality . In which case we have continuity in $0$ as desired.
\end{proof}


In the following we want to prove that a fractional Gaussian field has a $L^2$ density if and only if $s > \frac{d_h}{2d_w}$, but for this we need white noise.

\begin{defi}\label{whiteNoise}
	Let $(X, \mathcal{E}, \mu)$ be $\sigma$-finite measure space. A white noise based on $\nu$ is a random set function $W$ on the sets $A \in \mathcal{E}$ with $\nu(A) < \infty$ such
	\begin{enumerate}[(i)]
		\item $W(A)$ is a $N(0, \mu(A))$ random variable \label{whiteNoiseIsNormal}
		\item If $A \cap B = \emptyset$ then $W(A)$ and $W(B)$ are independent and $W(A \cap B) = W(A) + W(B)$ \label{whiteNoiseIsI.I.D}
	\end{enumerate}
\end{defi}

Now we will show that such a process exists on $(K, \mathcal{K}, \mu)$ where  $\mathcal{K}$ is the Borel $\sigma$-algebra on $K$. 

We will think of white noise as a Gaussian process indexed by the set $\mathcal{A} = \cbr{W(A) : A \in \mathcal{K}, \mu(A) < \infty}$. From \cref{whiteNoise}, (\ref{whiteNoiseIsNormal}) and (\ref{whiteNoiseIsI.I.D}) it is a mean-zero Gaussian process with covariance matrix
\begin{equation}
	\mathbf{\Sigma} = (\e \left[W(A) W(B)\right])_{W(A), W(B) \in \mathcal{A}} = (\mu(A \cap B))_{W(A), W(B) \in \mathcal{A}}.
\end{equation}
Now let $A_1, \ldots, A_n \in \mathcal{K}$ and $a_1, \ldots, a_n$ be real numbers then
\begin{equation}
	\sum_{i,j} a_i a_j \mathbf{\Sigma}_{i,j} = \sum_{i,j}a_ia_j \int_{K}\mathds{1}_{A_i} \cd \mathds{1}_{A_j} \dd \mu = \int_{K} \left(\sum_{i} a_i \mathds{1}_{A_i}\right)^2 \dd \mu \geq 0
\end{equation}
Thus $\mathbf{\Sigma}$ is positive definite, so there exists a probability space $(\Omega, \mathcal{F}, p)$ and a mean-zero Gaussian process $\cbr{W(A)}$ on $(\Omega, \mathcal{F}, p)$ such that $W$ satisfies \cref{whiteNoise} (\ref{whiteNoiseIsNormal}) and (\ref{whiteNoiseIsI.I.D}). 

This means that for $f \in L^{2}(K, \mathcal{K}, \mu)$ the integral with respect to $W$, $W(f) = \int_{K} f \dd W$ is a well defined Gaussian random variable. Moreover, the Gaussian measure $W$ gives rise to an isonormal Gaussian family $\cbr{W(f), f \in L^{2}(K, \mathcal{K}, \mu)}$ with the covariance function
\begin{equation}
	\e \left[W(f) W(g)\right] = \int_{K} fg \dd \mu
\end{equation}


In the following we want to define a fractional Brownian field, and we want to do it as the integral of the fractional Riesz kernel defined in \cref{fractionalRieszKernel}, thus we need $G_s(x, \cd ) \in L^{2}(K, \mu)$. 

So let $s > \frac{d_h}{2d_w}$ we choose $s$ as we do since later we need $\sum \lambda_j^{-s}$. We observe that from \cref{boundsOfRieszKernel} we have that $G_s(x, \cd ) \in L^{2}(K, \mu)$. Moreover we have that
	\begin{equation}
		G_s(x,y) = \frac{1}{\Gamma(s)}\int_{K} t^{s-1}p_t(x,y)\dd t = \frac{1}{\Gamma(s)}\int_{K} t^{s-1}\sum_{j = 1}^{\infty} e^{-\lambda_j t} \phi_j(x)\phi_j(y)\dd t = \sum_{j = 1}^{\infty} \frac{1}{\lambda_j^{-s}} \phi_j(x)\phi_j(y).
	\end{equation}
	recalling that $\frac{1}{\Gamma(s)}\int_{K} t^{s-1} e^{-\lambda_jt} \dd t = \lambda_j^{-s}$ and \cref{convergenceOfEigenvaluesOfLaplcianBounds}.

\begin{defi}
	Let $s > \frac{d_h}{2d_w}$. The fractional Brownian field with parameter $s$ is defined as 
	\begin{equation}
		\widetilde{X}_s(x) = \int_{K} G_s(x,y)\dd W (y)
	\end{equation}
\end{defi}

Note that since $s > \frac{d_h}{2d_w}$ we can write
\begin{equation}
	\widetilde{X}_s(x) = \sum_{i = 1}^{\infty} \lambda_i^{-s}\phi_i(x)W_i
\end{equation}
where $W_i = W(\phi_i)$ is an i.i.d sequence of Gaussian random variables with mean 0 and variance 1.

We will now show that one can construct a fractional Gaussian Field from the fractional Brownian field
\begin{prop}
	Let $s > \frac{d_h}{2d_w}$. Then the Gaussian random field defined by
	\begin{equation}
		X_s(f) = \int_{K}f(x) \widetilde{X}_{s}(x) \dd \mu (x), \quad f \in S(K)
	\end{equation}
	has the law of a fractional Gaussian field with parameter $s$. Moreover, if there exists a Gaussian field $\left(\widetilde{X}_{s}(x)\right)_{x \in K}$ on $K$ with
	\begin{equation}
		\e \left[\int_{K}\widetilde{Y}_S(x)^2 \dd \mu(x)\right] < \infty
	\end{equation}
	such that the Gaussian field defined by
	\begin{equation}
		Y_s(f) = \int_{K}f(x) \widetilde{Y}_{s}(x) \dd \mu (x), \quad f \in S(K)
	\end{equation}
	has the law of fractional Gaussian field with parameter $s$, then $s > \frac{d_h}{2d_w}$
\end{prop}
\begin{proof}
	Assume $s > \frac{d_h}{2d_w}$. The from Fubini's theorem we have for every $f \in S(K)$ that a.s 
	\begin{equation}
		\int_{K}f(x)\widetilde{X}_{s}(x) \dd \mu(x) =\int_{K}f(x) \int_{K} G_s(x,y) \dd W(y) \dd \mu(x) = \int_{K}(-\Delta)^{-s} f(y) \dd W(y)
	\end{equation}
	in which case by definition of $W$ we have that $X_s(f) = \int_{K}f(x)\widetilde{X}_{s}(x) \dd \mu(x)$ is a Gaussian random variable with mean zero and variance
	\begin{equation}
		\int_{K}\abs{(-\Delta)^{-s}f(x)}^{2}\dd \mu(x).
	\end{equation}
	and thus a fractional Gaussian field with parameter $s$. 
	
	\medskip
	Now assume that there is a Gaussian field $\left(\widetilde{X}_{s}(x)\right)_{x \in K}$ on $K$ with
	\begin{equation}
		\e \left[\int_{K}\widetilde{Y}_S(x)^2 \dd \mu(x)\right] < \infty
	\end{equation}
	such that the Gaussian field defined by
	\begin{equation}\label{theOneAbove}
		Y_s(f) = \int_{K}f(x) \widetilde{Y}_{s}(x) \dd \mu (x), \quad f \in S(K)
	\end{equation}
	has the law of fractional Gaussian field with parameter $s$. Then by spectral decomposition and equation \ref{theOneAbove} we have
	\begin{equation}
		\widetilde{Y}_s(x) = \sum_{i= 1}^{\infty} \inn{\widetilde{Y}_{s}}{\phi_i} \phi_i(x) = \sum_{i = 1}^{\infty} Y_s(\phi_i)\phi_i(x).
	\end{equation}
	
	Note that $(Y(\phi_i))_{i \geq 1}$ is a sequence of independent Gaussian random variables with mean zero and variance $(\lambda_{i}^{-2s})_{i \geq 1}$ since by \cref{fracGaussianField} and \cref{fracLaplacianProperties} we have that
	\begin{equation}
		\e \left[Y_s(\phi_i)\right] =\e \left[\int_{K} \phi_i \widetilde{Y}_s \dd \mu \right] = \e \left[\int_{K} \phi_i(x) \sum_{j = 1}^{\infty} \lambda_j^{-s}\phi_j(x)W_i \dd \mu(x) \right]= \e \left[W_i\right] =0
	\end{equation}
	and
	\begin{equation}
		\e \left[Y_s(\phi_i)Y_s(\phi_j)\right] = \int_{K}  (-\Delta)^{-s}\phi_i (-\Delta)^{-s}\phi_j \dd \mu = \frac{1}{(\lambda_{j}\lambda_i)^{s}} \int_{K} \phi_i \phi_j \dd \mu = \frac{1}{(\lambda_{j}\lambda_i)^{s}} \delta_{ij}.
	\end{equation}
	Since $\e \left[\int_{K} \widetilde{Y}_{s}(x)^2 \dd \mu (x)\right] < \infty$  we find that
	\begin{align}
		\e \left[\int_{K} \widetilde{Y}_{s}(x)^2 \dd \mu (x)\right]  &= \e \left[\int_{K} \left(\sum_{i = 1}^{\infty} Y_s(\phi_i)\phi_i(x)\right)^{2}\dd\mu(x)\right] \\
		& = \e \left[\sum_{i = 1}^{\infty} Y_s(\phi_i)^2\int_{K} \phi_i(x)\phi_i(x) \dd \mu (x) \right] \\
		& = \sum_{i = 1}^{\infty}\e \left[Y_s(\phi_i)^2\right] \\
		& = \sum_{i = 1}^{\infty} \frac{1}{\lambda_i^{-2s}} < \infty
	\end{align}
	where we use that $\int_{K} (\sum_{i = 1}^{\infty} Y_s(\phi_i)\phi_i(x))^2 \dd \mu(x) = \int_{K} \sum_{i = 1}^{\infty} Y_s(\phi_i)^2 \phi_i(x)^2 \dd \mu (x)$ which stems from the fact that $\int_{K}\phi_i \phi_j \dd \mu = \delta_{ij}$ and $Y_s(\phi_i)$ is a constant with respect to $\mu(x)$.
	
	We now have that
	\begin{equation}
		\e \left[\int_{K} \widetilde{Y}_{s}(x)^2 \dd \mu (x)\right] = \sum_{i = 1}^{\infty} \frac{1}{\lambda_i^{-2s}} < \infty
	\end{equation}
	in which case by \cref{convergenceOfEigenvaluesOfLaplacian} we have that $s > \frac{d_h}{2d_w}$
\end{proof}